\honest{The Futility of Emergence}{Eliezer Yudkowsky}{2007}

Phlogiston and vitalism are easy to judge in hindsight, so I dare to name a current theory with the same flaw: emergence. It is usually defined as the study of systems whose high-level behaviour arises from the interaction of many low-level elements, and I quote Wikipedia: ``The way complex systems and patterns arise out of a multiplicity of relatively simple interactions.'' Taken literally, that fits everything above the level of quarks. Pointing at a market crash and saying ``It's not a quark!'' explains nothing, and neither does ``It's an emergent phenomenon!'' \nb{It fits the one sentence quoted. The article's next section gave narrower definitions. On C. D. Broad's, from 1925, an emergent property ``cannot, even in theory, be deduced'' from the properties of its parts, and calling something emergent is a claim that can turn out false.}

I object to the noun, not the verb. ``X emerges from Y'' is fine if Y is a specific model, as gravity arises from the curvature of spacetime in general relativity, or chemistry from atoms in quantum electrodynamics. Saying that gravity depends on ``arisence'' would be no explanation. But people commonly use ``emergence'' as an explanation in itself. I have ``lost track'' of how often I have heard ``Intelligence is an emergent phenomenon!'' I quote no one. The phrase has every sign of a mysterious answer: it makes no new prediction about real minds, it satisfies curiosity without feeding it, it has no moving parts, those who say it ``confess their ignorance of the internals, and take pride in it,'' and intelligence stays as mysterious as before.

A fun exercise: delete ``emergent'' and see whether the sentence changes. ``Human intelligence is an emergent product of neurons firing'' becomes ``a product of neurons firing.'' Better still, say that we can predict some of an ant colony's behaviour from models of individual ants. \nb{That is close to Mark Bedau's 1997 definition of weak emergence: derivable from the parts, but only by running them forward.}

Another exercise: replace ``emergent'' with ``magical.'' ``Life is an emergent phenomenon'' becomes ``Life is a magical phenomenon,'' and I ask whether each pair does not ``fit exactly the same set of outcomes.'' \nb{Not in the sense Broad gave the word. Broad's ``Emergent Vitalism'' denied the special life-substance, or ``Entelechy,'' that substantial vitalism assumed, and a life-substance is the error the previous post charged vitalism with. The Wikipedia article the post links also quoted Bedau, a defender of weak emergence, saying that strong emergence ``is uncomfortably like magic.''}

Emergence is popular because it is ``the junk food of curiosity.'' I give no evidence for this. Humans are still humans even after a few science classes in college: freed from settled science, they do what their ancestors did, dressed in the literary genre of science.
